\documentclass[11pt,aspectratio=169]{beamer}

\usetheme{Madrid}
\usecolortheme{seahorse}

\usepackage[utf8]{inputenc}
\usepackage{amsmath,amssymb}
\usepackage{xcolor}
\usepackage{booktabs}
\usepackage{tikz}
\usetikzlibrary{arrows.meta,positioning,fit,backgrounds}

\setbeamertemplate{navigation symbols}{}
\setbeamertemplate{footline}[frame number]

\title{Beyond the Syllabus}
\subtitle{Multiplication and the FFT: A Solved Conjecture, an Open Problem, and a Dynamic
Closest Pair\\
\small Week 7 Supplement -- TOAE209/TOAH209: Design and Analysis of Algorithms}
\author{Foreign Trade University}
\date{}

\begin{document}

\begin{frame}
  \titlepage
\end{frame}

\begin{frame}{Outline}
  \tableofcontents
\end{frame}

% =================================================================
\section{Motivation}
% =================================================================

\begin{frame}{Why look beyond the syllabus?}
  \footnotesize
  \begin{itemize}
    \item This week you proved two $\Theta(n\log n)$ bounds: the \textbf{closest pair of
    points} via the strip argument, and \textbf{polynomial multiplication} via the FFT
    (evaluate at roots of unity, multiply pointwise, interpolate back).
    \item You also saw Karatsuba's $\Theta(n^{\log_2 3})$ integer multiplication --
    Week 1's algorithm, revisited as one point on a whole \emph{hierarchy} of
    multiplication algorithms.
    \item Three questions a curious student should ask next:
    \begin{enumerate}
      \item Karatsuba beat grade-school; the FFT-based approach beats Karatsuba. Is there
      an algorithm that beats \emph{that}, reaching $O(n\log n)$ for integers too?
      \item Is $\Theta(n\log n)$ actually \textbf{proven optimal} for the FFT itself, the
      way Week 3 proved Union-Find's bound optimal?
      \item Your closest-pair algorithm assumed a fixed, static point set. What if points
      keep arriving and leaving?
    \end{enumerate}
  \end{itemize}
  \begin{alertblock}{Note on scope}
    This deck deliberately does \emph{not} revisit the matrix-multiplication exponent
    $\omega$ (Strassen, AlphaTensor/AlphaEvolve) -- that story is already told in Week 2's
    supplement. Every claim below is sourced to a specific, checked paper.
  \end{alertblock}
\end{frame}

% =================================================================
\section{Closing the Loop: Integer Multiplication Reaches $O(n\log n)$}
% =================================================================

\begin{frame}{Recap: the hierarchy you have built so far}
  \small
  \begin{center}
  \begin{tabular}{@{}lll@{}}
    \toprule
    Algorithm & Time & Where it comes from \\
    \midrule
    Grade-school & $\Theta(n^2)$ & direct definition of multiplication \\
    Karatsuba (Week 1) & $\Theta(n^{\log_2 3}) \approx \Theta(n^{1.585})$ & 3 products
    instead of 4 \\
    Toom--Cook ($k$-way split) & $\Theta(n^{1+\varepsilon})$, any $\varepsilon>0$ & Karatsuba
    generalized to $k$ pieces \\
    Sch\"onhage--Strassen (1971) & $O(n\log n\log\log n)$ & FFT over integers (this week's
    idea) \\
    \textbf{Harvey--van der Hoeven} & $\boldsymbol{O(n\log n)}$ & next slide \\
    \bottomrule
  \end{tabular}
  \end{center}
  Each row is a strictly better asymptotic bound -- and, as you will see, each is
  optimized for a \emph{different} range of $n$ in practice.
\end{frame}

\begin{frame}{Schönhage \& Strassen's 1971 conjecture, confirmed in 2021}
  \begin{block}{Theorem (Harvey \& van der Hoeven, \emph{Annals of Mathematics} 193(2),
  2021, pp.\ 563--617)}
    Two $n$-bit integers can be multiplied using $O(n\log n)$ bit operations. \emph{DOI:}
    \texttt{10.4007/annals.2021.193.2.4}.
  \end{block}
  \begin{itemize}
    \item In 1971, Schönhage and Strassen introduced the $O(n\log n\log\log n)$ FFT-based
    algorithm you studied this week's ideas from, and \textbf{conjectured} $O(n\log n)$
    was the truth -- then it stood open for \textbf{48 years}.
    \item Harvey and van der Hoeven (2019, published 2021) finally constructed such an
    algorithm, confirming the conjecture.
  \end{itemize}
\end{frame}

\begin{frame}{The key idea: replace complex roots of unity with something exact}
  \begin{itemize}
    \item This week's FFT multiplies coefficient vectors by evaluating at complex
    $n$-th roots of unity $\omega_n = e^{2\pi i/n}$. Done with floating-point
    arithmetic, this accumulates \textbf{rounding error} that grows with $n$ --
    fine for polynomials with modest coefficients, fatal for multiplying huge
    exact integers.
    \item Schönhage--Strassen's fix: work in a ring $\mathbb{Z}/(2^m+1)\mathbb{Z}$ that
    has an exact analogue of a root of unity (a power of 2), so the transform is
    \textbf{exact} integer arithmetic, not floating point -- at the cost of an extra
    $\log\log n$ factor from recursively multiplying the (smaller) transformed
    coefficients.
    \item Harvey--van der Hoeven's contribution is a more intricate recursive choice of
    such rings (multi-dimensional FFTs over carefully engineered finite rings) that
    removes even that $\log\log n$ factor -- full technical machinery is well beyond this
    course, but the \emph{shape} of the idea -- ``do the FFT idea, but over an exact
    ring instead of $\mathbb{C}$'' -- is exactly this week's convolution trick, pushed
    one level further.
  \end{itemize}
\end{frame}

% =================================================================
\section{Theory vs.\ Practice: When Does the Optimal Algorithm Actually Win?}
% =================================================================

\begin{frame}{A cautionary crossover point}
  \begin{alertblock}{How large does $n$ have to be?}
    Harvey and van der Hoeven's algorithm only beats Schönhage--Strassen (and Fürer's
    2007 intermediate algorithm) once the integers being multiplied have more than
    \[
      2^{1729^{12}} \text{ bits.}
    \]
  \end{alertblock}
  \begin{itemize}
    \item For scale: the number of particles in the observable universe is estimated at
    about $2^{270}$. The crossover point is unimaginably larger than that.
    \item Source: D.\ Harvey, ``Integer multiplication in time $O(n\log n)$'' (talk
    slides, UNSW 2019); reported in \emph{Communications of the ACM}, ``Multiplication
    Hits the Speed Limit'' (Jan.\ 2020).
    \item Compare Week 3's $\alpha(n)$ table: there, a function that is \emph{unusably
    slow-growing in the good sense} (never exceeds 5 for any $n$ you could write down).
    Here, a crossover point that is \emph{unusably large in the bad sense}. Both are
    ``$O(1)$ for all practical purposes'' -- just in opposite directions.
  \end{itemize}
\end{frame}

\begin{frame}{What multiplication libraries actually run}
  \begin{itemize}
    \item No production library implements Harvey--van der Hoeven: the crossover point
    above puts it firmly in ``theoretical gold, not deployable'' territory.
    \item GMP (used by Python's \texttt{int}, SageMath, and most computer-algebra
    systems) instead switches, at \emph{measured} thresholds tuned per CPU, between
    grade-school $\to$ Karatsuba $\to$ Toom--Cook $\to$ FFT/Schönhage--Strassen-style
    methods as $n$ grows -- exactly the hierarchy from two slides ago, chosen
    empirically rather than by asymptotics alone.
    \item Recent (2024) work continues to refine this in new directions, e.g.\ evaluating
    NTT (Number-Theoretic Transform, an exact finite-field FFT)/Karatsuba \emph{hybrid}
    designs for $>16{,}000$-bit integers on FPGA hardware -- the same idea this week's FFT
    teaches you, retargeted at silicon.
  \end{itemize}
  \begin{exampleblock}{The lesson}
    ``Asymptotically optimal'' and ``fastest for the sizes you will ever multiply'' are
    different claims. This week's Karatsuba/FFT material is the one that actually wins in
    practice, for essentially every $n$ a computer will ever see.
  \end{exampleblock}
\end{frame}

% =================================================================
\section{Is the FFT Itself Optimal? A Genuinely Open Problem}
% =================================================================

\begin{frame}{Recall Week 3's move: prove a matching lower bound}
  \begin{itemize}
    \item For Union-Find, an upper bound ($O(m\,\alpha(n))$, Tarjan) was later matched by
    a lower bound (Fredman \& Saks, 1989, cell-probe model) -- so we \emph{know}
    Union-Find is optimal.
    \item Natural question: is $\Theta(n\log n)$ for computing the Discrete Fourier
    Transform of $n$ numbers similarly \textbf{proven optimal}?
  \end{itemize}
  \begin{alertblock}{Surprising answer: not in general}
    Whether \emph{every} algorithm computing the DFT needs $\Omega(n\log n)$ operations,
    in a fully general model of computation, is a genuinely \textbf{open problem} in
    complexity theory. Unlike Union-Find, this loop has not been closed.
  \end{alertblock}
\end{frame}

\begin{frame}{What \emph{is} known}
  \small
  Unconditional lower bounds exist only for \emph{restricted} models of computation:
  \begin{itemize}
    \item \textbf{Algebraic circuit model.} If you only count additions/multiplications
    in a linear circuit, $\Omega(n\log n)$ \emph{is} known and matches the FFT's upper
    bound -- so within that specific (very natural) model, the FFT is provably optimal.
    \item \textbf{Restricted unitary-gate model.} Ailon (2013), ``A Lower Bound for
    Fourier Transform Computation in a Linear Model Over $2{\times}2$ Unitary Gates
    Using Matrix Entropy'' (arXiv:1305.4745), proves $\Omega(n\log n)$ when each
    computation step is restricted to a $2{\times}2$ unitary operation on two
    coordinates at a time -- a model chosen to mimic how real FFT circuits are built.
    \item \textbf{Online, cell-probe model.} Clifford \& Jalsenius, ``Tight Cell-Probe
    Bounds for Online Integer Multiplication and Convolution'' (ICALP 2011,
    arXiv:1101.0768), prove a \emph{matching} lower bound specifically for the
    \emph{online} version (coefficients of one input revealed one at a time) --
    using the very same cell-probe technique as Fredman--Saks in Week 3.
  \end{itemize}
  No one has yet ruled out a fundamentally different, non-algebraic algorithm beating
  $n\log n$ for the general, offline DFT.
\end{frame}

% =================================================================
\section{Closest Pair Goes Parallel and Dynamic}
% =================================================================

\begin{frame}{Recap: this week's static, single-shot algorithm}
  \begin{itemize}
    \item This week's closest-pair algorithm assumes: (a) all $n$ points are known in
    advance, and (b) you compute the answer once, sequentially.
    \item Real applications -- ride-sharing dispatch, collision detection in a physics
    simulation, nearest-neighbor indexes -- need the closest pair \emph{as points keep
    being added and removed}, and ideally computed on many cores at once.
  \end{itemize}
\end{frame}

\begin{frame}{A parallel, batch-dynamic closest pair}
  \begin{block}{Wang, Yu, Gu \& Shun, ``A Parallel Batch-Dynamic Data Structure for the
  Closest Pair Problem'' (SoCG 2021; arXiv:2010.02379)}
    A data structure that maintains the closest pair of a point set under
    \textbf{batches} of insertions and deletions, with each batch update running in
    $O(\log^2 n)$ \emph{parallel depth} (span) using near-linear total work, rather than
    recomputing the static $O(n\log n)$ algorithm from scratch after every change.
  \end{block}
  \begin{itemize}
    \item The core reuse of this week's idea: the same ``strip near the dividing line''
    argument still bounds how much local structure a single update can disturb -- it
    is generalized, not discarded.
    \item ``Batch-dynamic'' means: update $b$ points at once in
    $O\!\left(\log^2 n + \frac{b}{p}\log^2 n\right)$-style parallel time with $p$
    processors, rather than paying a full sequential pass per point.
  \end{itemize}
\end{frame}

% =================================================================
\section{Summary}
% =================================================================

\begin{frame}{Summary}
  \footnotesize
  \begin{enumerate}
    \item \textbf{A 48-year conjecture, closed:} Harvey \& van der Hoeven (2021) built an
    exact, $O(n\log n)$ integer-multiplication algorithm, confirming Schönhage \&
    Strassen's 1971 guess -- by pushing this week's ``FFT over an exact ring instead of
    $\mathbb{C}$'' idea further.
    \item \textbf{Optimal in theory, irrelevant in practice (so far):} its crossover
    point, $2^{1729^{12}}$ bits, dwarfs any number a computer will ever multiply --
    real libraries still run tuned Karatsuba/Toom--Cook/Schönhage--Strassen hybrids.
    \item \textbf{An open loop:} unlike Union-Find (Week 3), no one has proven a general
    $\Omega(n\log n)$ lower bound for the DFT -- only for restricted models, using (among
    other tools) the same cell-probe technique from Week 3.
    \item \textbf{Static $\to$ dynamic \& parallel:} this week's closest-pair strip
    argument extends to a 2021 parallel batch-dynamic data structure, the same upgrade
    path Week 3's supplement traced for Union-Find/connectivity.
  \end{enumerate}
  \begin{alertblock}{Looking ahead}
    Next week (Kleinberg \& Tardos, Ch.~6) turns from divide-and-conquer's ``split,
    recurse, combine'' to \textbf{dynamic programming}'s ``solve every subproblem once,
    reuse the answers'' -- a different way of taming an exponential search space.
  \end{alertblock}
\end{frame}

\end{document}
